% ---------------------------------------------------------------
%  MNRAS version, converted from the REVTeX source by
%  Plotting_Scripts/to_mnras.py.  Do not edit this file by hand:
%  edit the REVTeX source and re-run the converter, or the two will
%  drift apart.
% ---------------------------------------------------------------
\documentclass[fleqn,usenatbib]{mnras}
\usepackage{graphicx}
\usepackage{amsmath,amssymb,bm}
\usepackage{txfonts}

\newcommand{\Esym}{E_{\mathrm{sym}}}
\newcommand{\Lsym}{L_{\mathrm{sym}}}
\newcommand{\DRnp}{\Delta R_{np}}
\newcommand{\Msun}{M_\odot}
\newcommand{\Teff}{T_{\mathrm{eff}}}
\newcommand{\fmc}{\mathrm{fm}^{-3}}

\title[Symmetry energy from NS cooling]{Constraining the nuclear symmetry energy from neutron star thermal evolution: the role of late-time heating, envelope composition, and extrapolation bias}

\author[Singh \& Nandi]{
Nisha Singh$^{1}$\thanks{E-mail: ns515@snu.edu.in}
and Rana Nandi$^{2}$
\\
$^{1}$Department of Physics, Shiv Nadar Institution of Eminence, Gautam Buddha Nagar, Uttar Pradesh 201314, India\\
$^{2}$Department of Physics, Indian Institute of Technology Delhi, Hauz Khas, New Delhi 110016, India
}

\begin{document}
\label{firstpage}
\maketitle

\begin{abstract}
We constrain the slope of the nuclear symmetry energy, $\Lsym$, by comparing
the observed thermal luminosities of sixteen isolated neutron stars with
cooling calculations built on relativistic mean field equations of state. The
isovector couplings $g_\rho$ and $\Lambda_{\omega\rho}$ are varied on a grid
in $(\Esym,\Lsym)$ at fixed isoscalar saturation properties for five
parametrizations, giving 864 equations of state that satisfy
$M_{\max}\ge 2\,\Msun$ and causality. We find
$\Lsym = 50$--$75$~MeV at 68\% credibility ($35$--$80$~MeV at 95\%) with
$\chi^2/\mathrm{dof}=1.60$, while $\Esym$ is only weakly constrained.

Reaching this required identifying three effects that dominate the result and
that we believe affect other cooling analyses. First, cooling curves terminate
when the star falls below the temperature floor of the microphysics tables, and
extrapolating past that point by holding the final temperature fixed biases the
comparison toward rapidly cooling models, because slowly cooling models survive
longer and are therefore scored against genuine rather than frozen predictions.
Second, without late-time heating no model reaches the ages of the oldest stars
in the sample, so those objects carry no information; including vortex creep
heating allows all sixteen stars to be used with no extrapolation. Third, two
of the hottest young objects are central compact objects with published carbon
atmospheres, and modelling them with iron envelopes accounts for a large part
of the residual. Correcting all three reduces $\chi^2/\mathrm{dof}$ from $4.70$
to $1.60$. We also report that the surface temperature adopted for
RX~J0822$-$4300 in earlier work corresponds to the hot component of a
two-component atmosphere fit rather than to the global surface.
\end{abstract}

\begin{keywords}
symmetry energy -- neutron star cooling -- direct Urca -- neutron skin
\end{keywords}


% =====================================================================
\section{Introduction}

The nuclear symmetry energy governs the energy cost of neutron--proton
asymmetry. About saturation density $\rho_0$ it is conventionally expanded as
\begin{equation}
S(\rho) = \Esym + \frac{\Lsym}{3}\,\frac{\rho-\rho_0}{\rho_0}
        + \frac{K_{\mathrm{sym}}}{18}\left(\frac{\rho-\rho_0}{\rho_0}\right)^2
        + \dots ,
\label{eq:expansion}
\end{equation}
with $\Esym \equiv S(\rho_0)$ and
$\Lsym \equiv 3\rho_0 (\partial S/\partial\rho)_{\rho_0}$. Because $\Lsym$
controls the pressure of neutron-rich matter just above saturation, it
simultaneously determines how far neutrons extend beyond protons in a heavy
nucleus \citep{Brown2000,Furnstahl2002,RocaMaza2011} and the proton fraction in
the core of a neutron star \citep{Lattimer1991}. Laboratory and astrophysical
determinations of $\Lsym$ have been reviewed extensively
\citep{LattimerLim2013,DanielewiczLee2014,Drischler2020,Essick2021}; the present
work adds an independent astrophysical route through thermal evolution.

The astrophysical lever exploited here is the direct Urca process
\citep{Lattimer1991},
\begin{equation}
n \rightarrow p + e^- + \bar\nu_e, \qquad p + e^- \rightarrow n + \nu_e ,
\end{equation}
which is kinematically permitted only when the degenerate Fermi momenta satisfy
the triangle inequality
\begin{equation}
k_{Fn} \le k_{Fp} + k_{Fe} .
\label{eq:durca}
\end{equation}
Equation~(\ref{eq:durca}) is the condition imposed in our calculation. It
reduces to the familiar $Y_p \ge 1/9$ only for $npe$ matter, where charge
neutrality gives $k_{Fp}=k_{Fe}$; once muons appear $k_{Fe} < k_{Fp}$ and the
threshold proton fraction rises. Our equations of state include muons, so the
$1/9$ figure should be read as a lower limit rather than the operative
threshold. Direct Urca exceeds modified Urca
\citep{FrimanMaxwell1979,Yakovlev2001} emissivity by several orders of
magnitude, so a star whose core crosses the threshold cools dramatically faster
\citep{YakovlevPethick2004,Page2004}. Since $Y_p$ at two to three times
saturation density is controlled largely by the symmetry energy \emph{slope},
the thermal history of the neutron star population is sensitive to $\Lsym$. The
closest published antecedent to the present analysis is
\citet{Sarkar2023}, which confronts fast cooling with the symmetry energy
in the light of PREX-2 \citep{PREX2}.

A corollary deserves emphasis: the threshold depends on how rapidly $S(\rho)$
\emph{grows} above saturation, not on its value at $\rho_0$. This is the
physical reason cooling constrains $\Lsym$ far better than $\Esym$, and we
regard the weak $\Esym$ sensitivity as a property of the observable rather than
a deficiency of the analysis.

% =====================================================================
\section{Equations of state}

\subsection{One Lagrangian, five parametrizations}

All five parametrizations (Table~\ref{tab:models}) share a single Lagrangian of
the standard non-linear relativistic mean field form
\citep{Walecka1974,BogutaBodmer1977,SerotWalecka1986} and differ only in their
couplings and in whether the quartic $\omega$ self-interaction $\zeta$
\citep{MullerSerot1996} is active:
\begin{align}
\mathcal{L} =\; & \bar\psi\Big[\gamma^\mu\big(i\partial_\mu - g_\omega\omega_\mu
 - \tfrac{1}{2} g_\rho \bm{\tau}\!\cdot\!\bm{\rho}_\mu\big)
 - (M - g_\sigma\sigma)\Big]\psi \nonumber\\
 & + \tfrac{1}{2}\big[(\partial\sigma)^2 - m_\sigma^2\sigma^2\big]
   - \tfrac{1}{3} b M (g_\sigma\sigma)^3 - \tfrac{1}{4} c (g_\sigma\sigma)^4
   \nonumber\\
 & - \tfrac{1}{4}\Omega_{\mu\nu}\Omega^{\mu\nu}
   + \tfrac{1}{2} m_\omega^2 \omega_\mu\omega^\mu
   + \tfrac{\zeta}{24}\big(g_\omega^2\omega_\mu\omega^\mu\big)^2 \nonumber\\
 & - \tfrac{1}{4}\bm{R}_{\mu\nu}\!\cdot\!\bm{R}^{\mu\nu}
   + \tfrac{1}{2} m_\rho^2 \bm{\rho}_\mu\!\cdot\!\bm{\rho}^\mu \nonumber\\
 & + \Lambda_{\omega\rho}\big(g_\omega^2\omega_\mu\omega^\mu\big)
     \big(g_\rho^2\bm{\rho}_\nu\!\cdot\!\bm{\rho}^\nu\big) ,
\end{align}
the last term being the $\omega$--$\rho$ cross coupling of
\citet{HorowitzPiekarewicz2001,ToddRutel2005}.
Writing $\Phi=g_\sigma\sigma$, $W=g_\omega\omega_0$ and $R=g_\rho\rho_{03}$,
the mean field equations in uniform matter are
\begin{align}
\frac{\Phi}{C_\sigma^2} + bM\Phi^2 + c\Phi^3 &= n_s ,
\label{eq:sigma}\\
\frac{W}{C_\omega^2} + 2\Lambda_{\omega\rho}R^2 W + \frac{\zeta}{6}W^3 &= n ,\\
\frac{R}{C_\rho^2} + 2\Lambda_{\omega\rho}W^2 R &= \frac{n_3}{2} ,
\end{align}
with $C_i^2 = (g_i/m_i)^2$, $n_3 = n_p-n_n$ and $M^\ast = M-\Phi$.

\begin{table}[t]
\caption{Isoscalar saturation properties of the five parametrizations. These are
held fixed throughout; only the two isovector couplings are varied.}
\label{tab:models}

\begin{tabular}{lccccl}
\hline
Set & $K$ (MeV) & $M^\ast/M$ & $\rho_0$ ($\fmc$) & $\zeta$ & Ref.\\
\hline
GM1 & 300.0 & 0.700 & 0.153 & 0 & \citep{GM1991}\\
GM2 & 300.0 & 0.780 & 0.153 & 0 & \citep{GM1991}\\
FSUGarnet & 229.6 & 0.578 & 0.153 & 0.0231 & \citep{ChenPiekarewicz2015}\\
IOPB-I & 222.6 & 0.595 & 0.149 & 0.0191 & \citep{IOPBI2018}\\
BigApple & 227.1 & 0.608 & 0.155 & 0.0009 & \citep{BigApple2020}\\
\hline
\end{tabular}

\end{table}

\subsection{Why only the isovector sector is varied}

The strategy of calibrating relativistic functionals against identical
isoscalar constraints while varying an isovector assumption follows
Chen and Piekarewicz \citep{ChenPiekarewicz2015}, who introduced one of the
parametrizations used here in exactly that way. Of the couplings in
Eq.~(\ref{eq:sigma}) and its partners, those in the isoscalar sector
($g_\sigma$, $g_\omega$, $b$, $c$, and $\zeta$ where present)
are fixed entirely by the published saturation properties $\rho_0$, $E/A$, $K$
and $M^\ast/M$ of each parametrization. Only the two isovector couplings
$g_\rho$ and $\Lambda_{\omega\rho}$ remain free, and they map one-to-one onto
$(\Esym,\Lsym)$. The scan is therefore well posed: at every grid point the
isoscalar sector is untouched and only the symmetry energy varies.

\subsection{Grid and selection}
\label{sec:grid}

For FSUGarnet, IOPB-I and BigApple, $\Esym$ is sampled at thirteen values from
$25$ to $43$~MeV in steps of $1.5$~MeV, and $\Lsym$ in steps of $5$~MeV. The
GM1 and GM2 sets were generated earlier on a finer and less regular $\Esym$
mesh, and retain six additional values ($\Esym = 26$, $30$, $32$, $36$, $38$
and $42$~MeV) together with two off-mesh $\Lsym$ points ($93.9$ and
$89.3$~MeV respectively). We have verified that this non-uniformity does not
affect the result: restricting the analysis to the common uniform sub-grid,
which removes $120$ of the $864$ equations of state, leaves the quoted $\Lsym$
interval unchanged at $[50,75]$~MeV and moves the $\Esym$ interval by half a
grid step. The full set is used below, and the sub-grid figure is reported in
Sec.~\ref{sec:systematics} as a robustness check.

The accessible range of $\Lsym$ is parametrization dependent, because below
some value no real $g_\rho$ exists: the lowest attainable $\Lsym$ is $35$~MeV
for BigApple, $45$~MeV for IOPB-I and $50$~MeV for FSUGarnet. Of the coupling
sets that yield a physical equation of state, \textbf{864} satisfy both
$M_{\max}\ge 2.0\,\Msun$ \citep{Demorest2010,Antoniadis2013,Fonseca2021} and
$v_s^2<1$ throughout the core, and it is these that enter the analysis. All
$864$ are scored on all sixteen stars.

\subsection{Crust and stellar structure}

The crust is the catalyzed Haensel--Zdunik--Dobaczewski \citep{HZD1989} /
Negele--Vautherin \citep{NegeleVautherin1973} table supplied with the cooling
code, whose highest-density row sits at $n = 0.0789~\fmc$. Fifteen of the
sixteen observed stars are isolated rather than accreting, so a catalyzed crust
\citep{BPS1971,DouchinHaensel2001} is the appropriate choice; this is data
driven rather than a default. Stellar structure follows from the
Tolman--Oppenheimer--Volkoff equations
\citep{Tolman1939,OppenheimerVolkoff1939}.

The limitation to record is different. A non-relativistic crust is joined to a
relativistic mean field core at a \emph{fixed} density $n=0.0789~\fmc$ for
every $(\Esym,\Lsym)$, whereas the true transition density falls as $\Lsym$
rises. The resulting inconsistency has been quantified by
Fortin \emph{et al.} \citep{Fortin2016} and we treat it as a systematic rather
than attempt a unified crust here. We have audited all 864 tables that enter the
analysis. Pressure and baryon density are monotonic in every one. Measuring the
\emph{fractional} pressure step across the crust--core boundary against the
fractional steps of its neighbours --- the scale-free comparison, since the
crust table is deliberately coarser than the core --- the join step is a median
$0.063$ times its neighbours, reaches $0.62$ at the 90th percentile, and never
exceeds $2.35$; no table exceeds three times. Causality is likewise satisfied:
no table has $v_s^2 > 1$ anywhere above
$\rho = 1.5\times 2.8\times10^{14}~\mathrm{g\,cm^{-3}}$, which is the density
above which the test is imposed and is well above the join.

% =====================================================================
\section{Thermal evolution and the observational comparison}

The observational sample is listed in Table~\ref{tab:obs}. Cooling curves are
computed for $M = 1.0$, $1.4$, $1.8$ and $2.0\,\Msun$ using
the one-dimensional general-relativistic thermal evolution code \textsc{NSCool}
\citep{Page2016}, with the neutrino emissivities of \citet{Yakovlev2001},
pairing suppression factors from \citet{Yakovlev1999}, and pair-breaking
emission as in \citet{Page2009}. The envelope follows
\citet{GPE1983,PCY1997}. The comparison statistic is
\begin{equation}
\chi^2 = \sum_{i=1}^{16} \min_{M}\left[
 \frac{\big(\log_{10}T^{\rm pred}_{\rm eff}(t_i,M)
 - \log_{10}T^{\rm obs}_{{\rm eff},i}\big)^2}{\sigma_i^2}
 + \frac{(M-1.4)^2}{\sigma_M^2}\right] ,
\label{eq:chi2}
\end{equation}
where each star selects the mass that fits it best but pays a prior penalty
for departing from $1.4\,\Msun$. Section~\ref{sec:mass} explains why the prior
is necessary.

Ages and surface temperatures for fourteen of the sixteen objects are taken
from the compilation of Potekhin \emph{et al.} \citep{Potekhin2020}. The two
remaining objects come from separate analyses: PSR~J0437$-$4715 from
\citet{Durant2012} and PSR~B0950$+$08 from \citet{Pavlov2017}.

\begin{figure}[t]
\includegraphics[width=\columnwidth]{figs/fig_cooling_curves.pdf}
\caption{Best-fit cooling curves at the four masses considered, together with
the sixteen observed stars. Red points mark the two central compact objects with
published carbon atmospheres, for which the carbon-envelope track (orange) is
used. All curves extend past the oldest star because late-time heating is
included; without it none reaches $\log_{10}t = 6.85$.}
\label{fig:curves}
\end{figure}

\begin{table}[t]
\caption{The sixteen stars used. Ages and temperatures for fourteen objects are
from the compilation of Potekhin \emph{et al.} \citep{Potekhin2020};
PSR\,J0437 \citep{Durant2012} and PSR\,B0950 \citep{Pavlov2017} are from separate
analyses. The last column gives the envelope composition adopted:
carbon for the two objects with published carbon atmospheres
\citep{HoHeinke2009,Klochkov2013}, iron otherwise.
The value for RX\,J0822 is the corrected whole-surface temperature discussed in
Sec.~\ref{sec:rxj}.}
\label{tab:obs}

\begin{tabular}{lcccc}
\hline
Source & $\log_{10}t$ & $\log_{10}T_{\rm eff}^\infty$ & $\sigma$ & env.\\
\hline
Cas\,A & 2.51 & 6.26 & 0.05 & C \\
XMMU\,J1732 & 3.48 & 6.25 & 0.10 & C \\
RX\,J0822 & 3.57 & 6.07 & 0.05 & Fe \\
1E\,1207 & 3.85 & 6.20 & 0.10 & Fe \\
PSR\,B0833 & 4.05 & 5.88 & 0.10 & Fe \\
PSR\,B1706 & 4.24 & 5.81 & 0.10 & Fe \\
PSR\,J0538 & 4.48 & 5.95 & 0.10 & Fe \\
PSR\,B2334 & 4.61 & 5.53 & 0.15 & Fe \\
PSR\,B0656 & 5.05 & 5.71 & 0.10 & Fe \\
Geminga & 5.53 & 5.75 & 0.10 & Fe \\
RX\,J1856 & 5.70 & 5.70 & 0.10 & Fe \\
PSR\,B1055 & 5.73 & 5.59 & 0.10 & Fe \\
RX\,J0720 & 5.78 & 5.72 & 0.10 & Fe \\
PSR\,J2043 & 5.88 & 5.22 & 0.15 & Fe \\
PSR\,J0437 & 6.85 & 5.34 & 0.15 & Fe \\
PSR\,B0950 & 7.24 & 5.02 & 0.20 & Fe \\
\hline
\end{tabular}

\end{table}

% =====================================================================
\section{Three effects that dominate the result}
\label{sec:three}

\subsection{Extrapolation past the end of a cooling curve}
\label{sec:clamp}

Cooling integrations terminate when the star drops below the temperature floor
of the microphysics tables, in our case $10^4$~K. Over the 864 equations of
state the median termination is $\log_{10}(t/\mathrm{yr}) = 6.23$ for direct
Urca with superfluidity and $6.54$ without it. The sample contains PSR~J0437$-$4715 at
$\log_{10} t = 6.85$ and PSR~B0950$+$08 at $7.24$. \emph{Without late-time
heating not one of the 864 equations of state reaches either age.}

If the analysis responds by returning the last computed temperature for all
later times, both stars are compared against a frozen number rather than a
prediction. This would be tolerable if it affected all models equally. It does
not. Slowly cooling scenarios remain above the floor longer and therefore run
to later times: over the 864 equations of state at all four masses, the
fraction of runs reaching $\log_{10}t = 6.85$ is $37.7\%$ for direct Urca
without pairing but \emph{exactly zero} for direct Urca with superfluidity. The slowly cooling models are
thus scored against genuine, and badly wrong, predictions while the rapidly
cooling ones receive a flattering frozen value. \textbf{Extrapolation of this
kind biases the comparison toward rapid cooling.}

We emphasise that lowering the temperature floor is not a remedy: below
$10^4$~K the conductivity, specific heat and emissivity tables do not extend,
and the integration fails within milliseconds of starting. The repair belongs
in the analysis, and the correct treatment is to decline to predict beyond the
computed range.

\subsection{Late-time heating}

The physical reason the oldest stars are unreachable is that no heating
mechanism was included. PSR~J0437$-$4715 ($7$~Myr, $2\times10^5$~K)
\citep{Durant2012} and PSR~B0950$+$08 ($17$~Myr, $10^5$~K) \citep{Pavlov2017} are
old and \emph{warm}, and remaining so requires late-time energy release
\citep{GonzalezReisenegger2010}. Including vortex creep heating
\citep{Alpar1984,ShibazakiLamb1989}, with the excess angular momentum as its
single parameter, extends every curve to $\log_{10} t \simeq 9.7$.
\textbf{All 864 equations of state then reach $\log_{10}t = 7.24$ for every
mass considered}, so all sixteen stars enter the fit with genuine predictions
and nothing is extrapolated. This also removes the mode-dependent coverage bias
of Sec.~\ref{sec:clamp} at a stroke.

\subsection{Envelope composition of the central compact objects}

Two of the hottest and youngest objects in the sample are central compact
objects \citep{DeLuca2017} with \emph{published carbon atmospheres}: the neutron
star in Cassiopeia~A \citep{HoHeinke2009} and XMMU~J173203.3$-$344518 in
HESS~J1731$-$347 \citep{Klochkov2013}. Modelling either
with an iron envelope is inconsistent with that spectroscopy, and both were
underpredicted by the iron-envelope calculation.

We therefore compute a second set of curves with a light-element envelope
\citep{PCY1997} for these two objects only. The light-element parameter is
\emph{not} fitted to the temperatures: it is calibrated so that the envelope
reproduces the surface enhancement that the Hernquist--Applegate relation
\citep{HernquistApplegate1984}, $T_{\rm eff}\propto (A^3/Z^4)^{1/4}$, predicts
for $^{12}$C relative to $^{56}$Fe, namely $+0.1351$~dex. The calibrated
value, $\eta = 6\times10^{-10}$, delivers a median surface enhancement of
$+0.138$~dex across the 864 equations of state; the value varies with the
cooling history of the individual track. The remaining fourteen stars retain iron envelopes;
allowing every star to choose its own envelope would amount to fitting sixteen
parameters to sixteen data points and is not done.

% =====================================================================
\section{The mass prior}
\label{sec:mass}

Equation~(\ref{eq:chi2}) profiles over mass, and it is worth stating plainly
what happens without the prior term. Free profiling assigns
$M = 1.0\,\Msun$ to eleven of the sixteen stars and yields
$\chi^2/\mathrm{dof} = 0.78$. Such a configuration is not credible: the
observed neutron star mass distribution peaks near $1.4\,\Msun$
\citep{OzelFreire2016} and $1.0\,\Msun$ stars are both rare and difficult to
form. The apparent quality of that fit is purchased with an unphysical mass
assignment.

With a Gaussian prior of width $\sigma_M = 0.3\,\Msun$ about $1.4\,\Msun$ the
selected masses become fourteen at $1.4$ and two at $1.8\,\Msun$, and
$\chi^2/\mathrm{dof} = 1.60$. For comparison, $\sigma_M = 0.2\,\Msun$ gives
$1.88$ and fixing all stars at $1.4\,\Msun$ gives $2.04$. We adopt
$\sigma_M = 0.3\,\Msun$ and quote $\chi^2/\mathrm{dof} = 1.60$ as our result,
regarding the free-profiling value as an unphysical lower bound.

% =====================================================================
\section{A correction to an adopted surface temperature}
\label{sec:rxj}

While checking the observational inputs against their sources we found that the
surface temperature previously adopted for RX~J0822$-$4300, the central compact
object in Puppis~A \citep{Becker2012}, was $\log_{10}T = 6.24$, corresponding to
$150$~eV. In the two-component atmosphere fit listed for this object in the
compilation of \citet{Potekhin2020}, from which that number derives, the
\emph{cold} component at $67$~eV is the one covering an emitting radius of
$12$~km, i.e.\ the whole star, while the $150$~eV component is a small hot spot.
% CHECK: add the original paper of the two-component fit as listed in
% Table 1 of Potekhin et al. (2020). It is NOT Zavlin et al. (1999), whose
% abstract gives a single H/He atmosphere fit with T = 1.6-1.9 MK. A fit with a
magnetized hydrogen atmosphere covering the entire surface gives instead
$T^\infty = 1.18\pm0.05$~MK, i.e.\ $\log_{10}T = 6.072$, which is also the
value adopted in \citet{Potekhin2020}.
% CHECK: cite the original nsmax fit here (analysis notes say
% "Shibanov et al. 2016", but the Shibanov et al. 2016 paper found,
% ApJ 831, 112, is about Calvera; take the reference from Potekhin 2020).
Other statistically acceptable fits span $0.79$--$1.34$~MK
($\log_{10}T = 5.90$--$6.13$), so the adopted $\sigma = 0.05$~dex is, if
anything, too tight for this object; this is listed among the systematics.

The effect is large. Re-running the final analysis with the old value
restored, and everything else unchanged, RX~J0822 becomes the single largest
contributor to the total $\chi^2$, at $15.9$ of $38.1$; with $6.072$ it
contributes $1.7$ of $22.3$. Correcting one input temperature therefore
accounts for most of the difference between those two totals. We recommend that
analyses of this object state explicitly which spectral component is being
used.

% =====================================================================
\section{Results}

\subsection{Fit quality}

Figure~\ref{fig:curves} shows the best-fit tracks against the data, and
Table~\ref{tab:progress} traces the effect of each correction. The total
improvement is a factor of $2.9$ in $\chi^2/\mathrm{dof}$, achieved entirely by
removing an extrapolation artifact, correcting one input temperature, and
adding physics that was absent, with no adjustable parameter tuned against the
temperatures.

\begin{table}[t]
\caption{Effect of each correction on fit quality. All entries use the same
864 equations of state.}
\label{tab:progress}

\begin{tabular}{lccc}
\hline
Treatment & stars used & $\chi^2/N$ & $\chi^2/\mathrm{dof}$ \\
\hline
extrapolated (frozen tail)      & 16, two frozen & 4.11 & 4.70 \\
no extrapolation                & 14             & 2.51 & 2.93 \\
$+$ RX~J0822 corrected          & 14             & 1.86 & 2.17 \\
$+$ late-time heating           & 16             & 1.47 & 1.68 \\
$+$ carbon envelopes, mass prior& 16             & 1.40 & 1.60 \\
\hline
\end{tabular}

\end{table}

\begin{figure}[t]
\includegraphics[width=\columnwidth]{figs/fig_chi2_vs_L.pdf}
\caption{Fit quality against $\Lsym$ for all 864 equations of state, separated
by parametrization. Dotted line: best fit. Dashed line: one unit of $\chi^2/N$
above it.}
\label{fig:chi2L}
\end{figure}

\begin{figure}[t]
\includegraphics[width=\columnwidth]{figs/fig_posterior_JL.pdf}
\caption{Relative likelihood over the scanned grid. The preference in $\Lsym$ is
pronounced; that in $\Esym$ is weak, as expected if the constraint originates in
the density dependence of $S(\rho)$ above saturation rather than in its value at
$\rho_0$.}
\label{fig:post}
\end{figure}

\subsection{Symmetry energy}

Because $\chi^2/\mathrm{dof}$ exceeds unity we inflate all uncertainties by the
Particle Data Group scale factor \citep{PDG2024}
$S = \sqrt{\chi^2/\mathrm{dof}} = 1.26$ before
forming credible intervals. Figures~\ref{fig:chi2L} and~\ref{fig:post} show the
fit quality and the resulting likelihood surface. Marginalising over all five
parametrizations with equal weight per parametrization, we obtain
\begin{align}
\Lsym &= 50\text{--}75~\mathrm{MeV} &(68\%),\nonumber\\
      &= 35\text{--}80~\mathrm{MeV} &(95\%),\\
\Esym &= 28.0\text{--}36.0~\mathrm{MeV} &(68\%),\nonumber\\
      &= 25.0\text{--}40.0~\mathrm{MeV} &(95\%).
\end{align}
The $\Esym$ interval spans most of the sampled range and should be read as a
weak constraint rather than a determination; as argued in the introduction,
this is expected. Intervals are quantised by the grid spacing, $5$~MeV in
$\Lsym$ and $1.5$~MeV in $\Esym$. The companion paper \citep{PaperII} quotes
$\Lsym = 45$--$70$~MeV for the cooling constraint alone, one grid step lower,
because it must restrict the marginalisation to the three parametrizations with
published meson masses and to $\Esym \ge 26$~MeV; the reason is set out there.

Our $\Lsym$ interval is consistent with, and comparable in width to,
determinations from nuclear structure: $35$--$70$~MeV from isobaric analog
states \citep{DanielewiczLee2014}, $40.5$--$61.9$~MeV from the concordance of
experimental, theoretical and observational analyses assembled in
\citet{LattimerLim2013}, and
$L \sim 60\pm25$~MeV as summarised in \citet{RocaMaza2011}. It is lower
than the value inferred from PREX-2 alone, $106\pm37$~MeV
\citep{Reed2021}, and compatible with the model-light astrophysical inference of
\citet{Essick2021} and the chiral effective field theory range of
\citet{Drischler2020}. We regard the agreement as a cross-check of the
cooling method rather than an improvement in precision.

A recent Bayesian analysis combining nuclear structure data and
INDRA--FAZIA isospin-diffusion measurements with chiral effective field theory
and astrophysical constraints \citep{Montefusco2026} bears directly on this
result. It finds that nucleonic direct Urca at or below $1.4\,\Msun$ has a
posterior probability below $1\%$, with $M_{\rm DU}=1.89^{+0.24}_{-0.21}\,\Msun$,
in agreement with the threshold masses of Sec.~\ref{sec:mdu}. Its symmetry
energy, $\Esym=29.4\pm0.7$~MeV, lies inside our interval, but its slope,
$\Lsym=29\pm8$~MeV, is lower than ours; the two 68\% intervals do not overlap
and meet only at our 95\% edge. Part of the difference is structural. The
lowest $\Lsym$ with a real $g_\rho$ is $35$, $45$ and $50$~MeV for BigApple,
IOPB-I and FSUGarnet, and our grid stops at $30$~MeV, so the region they favour
is sampled here only by GM1 and GM2. Moreover, with $\omega$--$\rho$ mixing the
slope at saturation and $S(n)$ at two to three times saturation, which is what
cooling measures, are tied together more rigidly than in their metamodel. Their
isospin-diffusion constraint alone, $\Lsym=41^{+14}_{-15}$~MeV, does overlap our
interval. Extending the scan to lower $\Lsym$ is the natural test of this
tension.

\begin{figure}[t]
\includegraphics[width=\columnwidth]{figs/fig_R14.pdf}
\caption{$R(1.4\,\Msun)$ under a flat prior over the surviving equations of
state and after the cooling constraint, compared with the NICER inference
\citep{Miller2021}.}
\label{fig:r14}
\end{figure}

\subsection{A note on the best-fitting parametrization}

The single best-fitting equation of state is GM2 at $\Esym = 32.5$~MeV,
$\Lsym = 80$~MeV. We draw no conclusion from the identity of that one point.
GM1 and GM2 are fitted to uniform nuclear matter alone and publish only the
ratios $g_i/m_i$ \citep{GM1991}; as a consequence they cannot be used for
finite-nucleus observables at all, for reasons set out in the companion
paper \citep{PaperII}. They are retained here because the cooling calculation
requires only the equation of state, for which they are complete, but any
statement connecting our cooling constraint to nuclear structure necessarily
rests on the three parametrizations that are complete in that stronger sense.

\subsection{Direct Urca threshold masses}
\label{sec:mdu}

For every equation of state the direct Urca threshold density $n_{\rm DU}$
follows from Eq.~(\ref{eq:durca}), and the threshold mass $M_{\rm DU}$ is the
mass of the star whose central density equals $n_{\rm DU}$, read off the
stable branch of the mass--radius curve. Within the 68\% region
($\Lsym=50$--$75$~MeV, $\Esym=28$--$36$~MeV) the median is
$M_{\rm DU}=1.77\,\Msun$ and $9.9\%$ of the equations of state open direct Urca
at or below $1.4\,\Msun$; for $35$ of the $152$ no stable star reaches the
threshold at all, all but one of them GM1 or GM2. The best-fitting equation of
state has $M_{\rm DU}=1.47\,\Msun$. This is the picture behind the fit: stars
near $1.4\,\Msun$ cool without direct Urca, while the two coldest objects,
which select $1.8\,\Msun$, lie above threshold.
% CHECK: replace 9.9% (fraction of grid points) by the posterior-weighted
% probability P(M_DU <= 1.4) from cooling_chi2_final.csv.

\subsection{Implied radius}

Figure~\ref{fig:r14} shows the result of pushing the posterior onto stellar
structure. It gives
$R(1.4\,\Msun) = 12.96$--$13.63$~km at 68\% ($13.20\pm0.39$~km), against a
flat-prior range of $12.18$--$14.49$~km, so the thermal data do constrain the
radius. This sits $1.0\sigma$ above the NICER inference
$12.45\pm0.65$~km \citep{Miller2021,Riley2021}, which we regard as agreement.

The comparison is meaningful rather than vacuous: $48$ of the $864$ equations of
state do produce $R(1.4) < 12.45$~km, all of them GM2 at
$\Lsym = 30$--$55$~MeV, and the best $\chi^2/N$ among them is $1.78$ against
$1.40$ globally. Small radii are therefore accessible to the model family and
are mildly disfavoured by the thermal data, rather than being excluded by
construction. We also note that $R(1.4\,\Msun)$ and $\Lsym$ are correlated at
$0.892$ across the grid, so the radius statement is largely the $\Lsym$
constraint in other units and is not independent information.

We do not quote a constraint on $M_{\max}$. The posterior mean
$2.23\pm0.21\,\Msun$ sits inside the flat-prior range
$[2.00,2.63]\,\Msun$, whose lower edge is set by our own selection cut.
Thermal evolution probes the direct Urca threshold at two to three times
saturation density, not the much higher central densities that set the maximum
mass.

% =====================================================================
\section{Systematic uncertainties}
\label{sec:systematics}

\paragraph{Isoscalar calibration of the $\zeta\neq0$ sets.} For FSUGarnet,
IOPB-I and BigApple the cubic and quartic $\sigma$ couplings and $\zeta$ were
obtained from the published saturation point with $C_\sigma^2$ and $C_\omega^2$
held at their published values. The resulting $\zeta$ differ slightly from the
published ones, and the incompressibility is reproduced to $1.4\%$ rather than
exactly. Refitting exactly to $(\rho_0,E/A,K,M^\ast/M)$ changes $M_{\max}$ by
at most $0.001\,\Msun$. Adopting the published $\zeta$ reproduces the published
maximum masses to $0.2\%$ (IOPB-I: $2.149\,\Msun$), raises the IOPB-I maximum
mass by $0.03\,\Msun$, and moves $M_{\rm DU}$ by at most $0.02\,\Msun$, well
below the resolution of the cooling comparison.
% CHECK: confirm the published zeta (FSUGarnet 0.0235, IOPB-I zeta_0=3.103,
% BigApple 0.0007) against the original parameter tables.

\paragraph{Temperature of RX~J0822$-$4300.} The spread of acceptable
whole-surface fits, $\log_{10}T=5.90$--$6.13$, is wider than the adopted
$\sigma=0.05$~dex.
% CHECK: rerun with sigma = 0.12 for RX J0822 and quote the change.

\paragraph{Crust consistency.} A fixed crust--core join density for all
$(\Esym,\Lsym)$, and a non-relativistic crust on a relativistic core; see
\citet{Fortin2016} for the size of this effect.

\paragraph{Grid non-uniformity.} As described in Sec.~\ref{sec:grid}, the GM1
and GM2 sets carry a finer $\Esym$ mesh than the other three. Repeating the
entire analysis on the common uniform sub-grid ($742$ equations of state)
returns $\Lsym = [50,75]$~MeV and $\Esym = [28.0,35.5]$~MeV, against
$[50,75]$ and $[28.0,36.0]$ for the full set. The constraint is therefore not
an artifact of the sampling.

\paragraph{Age uncertainties.} The statistic in Eq.~(\ref{eq:chi2}) carries a
temperature uncertainty per star but no age uncertainty. Propagating an assumed
age error through the local slope of the cooling curve, a uniform
$\sigma_{\log t}=0.3$ reduces the best $\chi^2/N$ from $3.64$ to $2.51$ in the
un-heated analysis, and even an implausibly large $0.5$ leaves it above $2$.
Age errors therefore matter but cannot by themselves account for the misfit.

\paragraph{Remaining input temperatures.} Five further entries in our adopted
compilation differ from the source tabulation \citep{Potekhin2020} by more than
$0.1$~dex, four of them colder than the values we used. Because our models
already run cold at late times, revising these would likely worsen rather than
improve the fit. We have not revised them here, as each requires a judgement
about whether the appropriate number is a blackbody fit, an atmosphere fit, or
a whole-surface model.

\paragraph{Heating parameter.} The vortex creep amplitude is a single global
parameter, not fitted per star.

\paragraph{Neutron $^3P_2$ pairing gap.} Repeating the entire analysis with an
alternative $^3P_2$ gap changes the best $\chi^2$ by $0.21$, i.e.\ $0.4\sigma$
after error inflation, and selects the same best-fitting equation of state. The
constraint is therefore insensitive to this choice within the pair of gaps
tested. For the wider uncertainty in the gaps see
\citet{Shternin2011,Ho2015}.

\paragraph{Cooling mechanism.} Profiling over the physics mode rather than
fixing it, the five scenarios considered span $\chi^2/N$ from $1.86$ to $2.60$
in the un-extrapolated analysis, a spread of $2.35\sigma$ after error
inflation. We therefore make no claim that the data select a cooling
mechanism, and note that any such claim drawn from an extrapolated analysis
would inherit the bias of Sec.~\ref{sec:clamp}.

\paragraph{Quality of the cooling runs.} All $3456$ runs entering the final
analysis terminated normally, every one reaching
$\log_{10}t \ge 9.58$ with a median of $9.68$, well beyond the oldest star. For
comparison, $1.0\%$ of the un-heated runs over the same equations of state
terminated abnormally; those runs are not used here.

% =====================================================================
\section{Conclusions}

Comparing 864 relativistic mean field equations of state with the thermal
luminosities of sixteen isolated neutron stars, and treating the comparison
carefully, we find $\Lsym = 50$--$75$~MeV at 68\% credibility
($35$--$80$~MeV at 95\%) with $\chi^2/\mathrm{dof}=1.60$, while $\Esym$ is only
weakly constrained.

We believe the methodological findings are of wider relevance. Extrapolating a
cooling curve beyond the point at which the integration stopped is not a
neutral choice: because slowly cooling models survive longer, it biases the
comparison toward rapid cooling, and in our case it inverted both the preferred
$\Esym$ and the ranking of cooling scenarios. Any analysis comparing models
with different termination times inherits the same bias. Including late-time
heating removes the problem at its source by making the oldest stars reachable.
Finally, objects with independently established atmospheric compositions should
be modelled with them; two such objects accounted for a large part of our
residual.

The neutron skins implied by these same parametrizations, and their confrontation
with PREX-2 and CREX, are the subject of the companion paper \citep{PaperII}.


\section*{Acknowledgements}

N.S. thanks Kenji Nishiwaki for supervision and support, and acknowledges Shiv
Nadar Institution of Eminence for a research fellowship. We thank Dany Page for
making \textsc{NSCool} publicly available. The numerical calculations were
performed on institutional computing resources at Shiv Nadar Institution of
Eminence.

\section*{Data availability}

The equation-of-state tables, cooling outputs and analysis products that
underlie this work are available from the authors on reasonable request.
The thermal evolution code is publicly available from its author and is
not redistributed here.

\bibliographystyle{mnras}
\bibliography{refs}

\label{lastpage}
\end{document}
